\subsection{Models and Distributed Learning}

\begin{frame}{Baseline: SGD Ridge sets the linear reference}

\begin{columns}[T,onlytextwidth]
	\column{.52\textwidth}
	\textbf{One global linear model}
	\[
		\widehat y_i = \mathbf{w}^{\mathsf T}\mathbf{x}_i+b
	\]
	\[
		J(\mathbf{w},b)=
		\frac{1}{N}\sum_{i=1}^{N}(\widehat y_i-y_i)^2
		+\lambda\lVert\mathbf{w}\rVert_2^2
	\]
	\[
		\mathbf{w}\leftarrow
		\mathbf{w}-\eta\nabla_{\mathbf{w}}J
	\]

	\column{.43\textwidth}
	\textbf{Why it is the baseline}
	\begin{itemize}
		\setlength{\itemsep}{4pt}
		\item The squared error directly penalizes year error.
		\item L2 shrinkage stabilizes correlated audio features.
		\item The objective is convex, so SGD has a clear target.
		\item Coefficients give an interpretable linear reference.
	\end{itemize}
\end{columns}

\medskip

\begin{block}{Limitation}
One hyperplane cannot express thresholds or feature interactions: the effect of a
feature is fixed for every song and every era.~{\scriptsize\autocite{hoerl1970ridge}}
\end{block}

\end{frame}

\begin{frame}{LightGBM: nonlinear boosting improves the baseline}

\small

\begin{columns}[T,onlytextwidth]
	\column{.52\textwidth}
	\textbf{Build an additive tree model}
	\[
		F_M(\mathbf{x})=F_0(\mathbf{x})
		+\sum_{m=1}^{M}\eta f_m(\mathbf{x})
	\]
	Each tree $f_m$ follows the negative loss gradient left by the
	previous trees.

	\medskip
	\textbf{Robust Huber objective}
	\[
	\rho_\delta(r)=
	\begin{cases}
		\frac{1}{2}r^2,& |r|\leq\delta,\\
		\delta(|r|-\frac{1}{2}\delta),& |r|>\delta.
	\end{cases}
	\]

	\column{.43\textwidth}
	\textbf{Advantages over Ridge}
	\begin{itemize}
		\setlength{\itemsep}{2pt}
		\item Tree splits learn nonlinear thresholds.
		\item A path through a tree represents feature interactions.
		\item Huber loss limits the influence of extreme year errors.
		\item Histogram splits reduce the cost of searching continuous values.
	\end{itemize}
\end{columns}

\medskip
\begin{block}{Practical effect}
Trees learn conditional rules such as ``loudness matters only at a certain
tempo''; Ridge can only add two fixed linear effects.
~{\scriptsize\autocite{ke2017lightgbm,huber1964robust}}
\end{block}

\end{frame}

\begin{frame}[shrink=12]{Ordinal-MoE: order plus era-specific experts}

\small

\begin{columns}[T,onlytextwidth]
	\column{.53\textwidth}
	\textbf{CORAL ordinal head: respect year order}
	\[
	\begin{aligned}
	p_k&=\sigma(s-\theta_k),\\
	\widehat y_{\rm ord}&=1922+\sum_{k=0}^{88}p_k
	\end{aligned}
	\]
	Here $p_k$ is the probability of passing ordered year threshold $k$.

	\medskip
	\textbf{Mixture of decade experts}
	\[
	\begin{aligned}
	\mathbf{g}&=\operatorname{softmax}(W_g\mathbf{x}+\mathbf{b}_g),\\
	\widehat y_{\rm moe}&=\sum_j g_j e_j(\mathbf{x})
	\end{aligned}
	\]
	The gate $g_j$ chooses how strongly each decade expert contributes.

	\column{.42\textwidth}
	\textbf{Why go beyond one tree ensemble?}
	\begin{itemize}
		\setlength{\itemsep}{2pt}
		\item Ordered thresholds prevent treating years as unrelated labels.
		\item Experts learn different feature--year relations by era.
		\item Soft gates allow gradual transitions between decades.
		\item Bounded expert outputs avoid implausible years.
	\end{itemize}
\end{columns}

\medskip
\[
\begin{aligned}
\widehat y&=0.20\,\widehat y_{\rm ord}+0.80\,\widehat y_{\rm moe},\\
\mathcal L&=.35\mathcal L_{\rm ord}+.45\rho_\delta\\
&\quad +.12\mathcal L_{\rm decade}+\text{auxiliary terms}.
\end{aligned}
\]
{\small Multiple supervised heads combine global order with local era
specialization.~{\scriptsize\autocite{cao2020coral,jacobs1991experts}}}

\end{frame}

\begin{frame}{Frozen test: three models, two kinds of improvement}

\begin{center}
\setlength{\tabcolsep}{12pt}
\begin{tabular}{l r r}
	\textbf{Model} & \textbf{MAE (years)} & \textbf{RMSE (years)} \\
	\hline
	SGD Ridge & $6.7465$ & $9.4074$ \\
	LightGBM & $5.8339$ & $\mathbf{8.3633}$ \\
	Ordinal-MoE blend & $\mathbf{5.4745}$ & $8.8143$ \\
\end{tabular}
\end{center}

\medskip

\begin{columns}[T,onlytextwidth]
	\column{.47\textwidth}
	\begin{block}{Lowest average absolute error}
	Ordinal-MoE reduces Ridge MAE by $18.9\%$. Its ordered, era-specific
	structure improves the typical prediction.
	\end{block}

	\column{.47\textwidth}
	\begin{block}{Lowest large-error penalty}
	LightGBM reduces Ridge RMSE by $11.1\%$ and has the best RMSE, indicating
	better control of the largest errors.
	\end{block}
\end{columns}

\medskip
{\small Same frozen artist-disjoint test split: $49{,}436$ tracks. Model
selection and preprocessing use no test labels.}

\end{frame}

\begin{frame}{Spark makes the model comparison scale}

\small

\begin{center}
\textbf{Parquet} $\longrightarrow$
\textbf{Spark partitions} $\longrightarrow$
\textbf{local computation} $\longrightarrow$
\textbf{aggregate/update} $\longrightarrow$
\textbf{distributed metrics}
\end{center}

\[
\nabla J =
\frac{1}{N}\sum_{p=1}^{P}
\underbrace{\sum_{i\in\mathcal P_p}\nabla\ell_i}_{\text{computed on partition }p}
+\lambda\mathbf{w}
\]

\begin{columns}[T,onlytextwidth]
	\small
	\column{.30\textwidth}
	\textbf{SGD Ridge}
	\begin{itemize}
		\setlength{\itemsep}{2pt}
		\item Local residual and gradient sums.
		\item Reduce before one global update.
	\end{itemize}

	\column{.30\textwidth}
	\textbf{LightGBM}
	\begin{itemize}
		\setlength{\itemsep}{2pt}
		\item SynapseML coordinates workers.
		\item Distributed split histograms.
	\end{itemize}

	\column{.34\textwidth}
	\textbf{Ordinal-MoE}
	\begin{itemize}
		\setlength{\itemsep}{2pt}
		\item Broadcast current head parameters.
		\item Reduce analytic mini-batch gradients.
	\end{itemize}
\end{columns}

\smallskip
\begin{block}{Why Spark matters}
The $515$k-track feature matrix stays partitioned instead of being collected
to one driver; caching and parallel computation support repeated training and
consistent evaluation.~{\scriptsize\autocite{zaharia2012rdd}}
\end{block}

\end{frame}
